% "You cannot rank what you can make" --- lifted verbatim from
% results/experimental.tex on 2026-10-25, where it sat as two subsections of
% the wet-lab section. Not one word was changed in the move.
%
% It is here because the restart made it load-bearing rather than supporting.
% The introduction's claim is that millions of candidates pass every filter and
% nothing ranks them; this section is the measurement behind that sentence, so
% it now comes before the library rather than after it.
\section{You cannot rank what you can make}
\label{sec:res:ranking}

\plan{\textbf{Section is yours to frame.} What is below is the evidence,
already written. What it is being asked to do is new: it has to carry the
intro's ``nearly impossible to rank which of those will be most functional''
from assertion to measurement, and hand off to the cost section
(Section~\ref{sec:res:cost}) --- if you \emph{could} rank them you would not
need a library, and if ranking fails the library is the only remaining move.

\textbf{The figure that was missing now exists.} The cross-enzyme comparison
had no figure script and was an inline table; it is drawn as
Figure~\ref{fig:transfer-scatter} below, from
\texttt{paper/figures/campaign\_panels.py}.}

\begin{figure}[t]
  \centering
  \includegraphics[width=0.78\textwidth]{figures/pdf/panel_c10.pdf}
  \caption{Each point is one metric, plotted by how well it predicted activity
  on PETase against how well it predicted on a second de novo serine hydrolase.
  On the diagonal a metric transfers; in the off-diagonal quadrants it points
  the wrong way. The \num{49} shared AF3 concepts anti-correlate at
  $\rho$-of-$\rho$s $= -0.604$, with 10 agreeing in sign against ${\sim}24$
  expected by chance. The \num{66} Rosetta concepts fail more mildly
  ($+0.035$, 33 of 66): no transferable signal rather than an inverted one.}
  \label{fig:transfer-scatter}
\end{figure}

\subsection*{Zero-shot prediction is lightly predictive}

\begin{figure}[t]
  \centering
  \includegraphics[width=\textwidth]{figures/pdf/petase_zeroshot.pdf}
  \caption{\textbf{Structure-prediction confidence against measured activity
  for 3{,}135 recombined PETase designs.}
  (a) The strongest zero-shot signal is local and state-specific: per-atom
  pLDDT over the active site of the \emph{tetrahedral adduct}. Selecting
  the top 5\% by this metric returns top-decile actives at $2.8\times$ the base
  rate. (b) Activity by design round, on the single screen run. (c) The
  control: distance from the design parent ranks activity better than any
  prediction does. The 20 open red points are a selected set re-assayed in a
  different run of the same assay and are never pooled into a correlation.}
  \label{fig:petase-zeroshot}
\end{figure}

Every design was folded by AF3 in two states --- the Michaelis complex and a
covalent \emph{tetrahedral adduct} --- and scored on 52 confidence and geometry
metrics. The adduct is not an acyl-enzyme intermediate, although the column is
named \texttt{acyl\_*}: AF3 strips the leaving hydrogens but keeps the leaving
oxygens, so the acid --OH survives and the bonded carbon stays
4-coordinate. Both enzymes in the transfer comparison below model the same
species, so that comparison is unaffected --- but the state is the adduct, and
calling it the acyl-enzyme intermediate would misstate which chemistry the
predictor was scored on. Table~\ref{tab:petase-zeroshot} gives the ranking against measured
activity.

\begin{table}[h]
  \centering
  \small
  \begin{tabular}{llrr}
    \toprule
    Predictor & What it measures & Spearman $\rho$ & $n$ \\
    \midrule
    \multicolumn{4}{l}{\emph{Zero-shot, AF3}} \\
    \quad acyl His107-NE2 pLDDT      & catalytic atom, acyl state & $+0.60$ & 2{,}548 \\
    \quad acyl active-site pLDDT     & active-site shell, acyl state & $+0.58$ & 2{,}548 \\
    \quad acyl PAE Ser66$\leftrightarrow$His107 & triad geometry & $-0.47$ & 2{,}548 \\
    \quad substrate interface PAE (min) & protein--ligand placement & $-0.43$ & 2{,}616 \\
    \quad substrate pTM               & global confidence & $+0.36$ & 2{,}616 \\
    \quad substrate ipTM              & global interface confidence & $+0.32$ & 2{,}616 \\
    \midrule
    \multicolumn{4}{l}{\emph{Zero-shot, sequence only}} \\
    \quad mutations from parent      & distance, no structure at all & $-0.65$ & 3{,}135 \\
    \midrule
    \multicolumn{4}{l}{\emph{Supervised, held out} (RidgeCV, 200 stratified splits)} \\
    \quad AF3 features               & & $+0.73 \pm 0.02$ & \\
    \quad Rosetta features           & & $+0.75 \pm 0.01$ & \\
    \quad AF3 $+$ Rosetta            & & $+0.80 \pm 0.01$ & \\
    \bottomrule
  \end{tabular}
  \caption{\textbf{Ranking power against measured activity.} Zero-shot rows are
  univariate Spearman on the screen run. Global confidence is roughly half as
  informative as a local, state-specific measurement at the catalytic atoms.}
  \label{tab:petase-zeroshot}
\end{table}

Three readings, in the order that matters.

\textbf{The useful signal is local and state-specific, not global.} pTM and ipTM
--- the metrics a design pipeline reaches for by default --- reach $\rho \approx
0.35$. Per-atom pLDDT at the catalytic histidine of the \emph{tetrahedral
adduct} reaches $0.60$. The gain comes from scoring the chemically
relevant state at the chemically relevant atoms, which is an argument for
building the covalent intermediate into the oracle rather than for folding
harder.

\textbf{It is a real signal, and it is light.} Selecting the top 5\% of the
library by active-site pLDDT returns a top-decile-active design 28\% of the
time against a 10\% base rate: a $2.8\times$ enrichment, falling to $1.9\times$
if 25\% of the library is kept. That is worth having when the alternative is
random ordering, and it is nowhere near a filter one would trust to replace an
assay.

\textbf{And a cheaper predictor beats it.} Mutation count from the parent, which
requires no structure prediction at all, ranks activity at $\rho = -0.65$ ---
better than any AF3 metric in the table (Figure~\ref{fig:petase-zeroshot}c).
This is the comparison a zero-shot claim has to clear and it does not clear it.
What survives is that the AF3 signal is not merely that confound re-expressed:
restricted to the 2{,}539 designs at least 20 mutations from the parent, active-site
pLDDT still gives $\rho = +0.59$. Within single 5-mutation bands it ranges from
$+0.05$ nearest the parent to $+0.59$ at \numrange{30}{34} mutations, and is
indistinguishable from zero in the three most distant bands
(\numrange{-0.01}{+0.09} beyond 50). Prediction adds something to distance; it
does not dominate it --- but it adds it only in the middle of the range, and the
band-wise view is what shows that the pooled $+0.59$ is not uniform.

\subsection*{\dots and it does not transfer to a second enzyme}

The $\rho = 0.60$ above is a property of this scaffold, not of the method. The
same two-state fold and the same metric pipeline were run against a second,
unrelated serine hydrolase --- \texttt{super}, a 160-aa de novo design
(Ser128/His95/Asp53, 4-methylumbelliferyl butyrate, 2{,}509 purified designs
measured over 2.51 decades). Correlating each metric's \emph{usefulness} on one
enzyme against its usefulness on the other asks whether a predictor chosen once
can be reused. It cannot.

\begin{table}[h]
  \centering\small
  \begin{tabular}{@{}llrrc@{}}
    \toprule
    State & Concept & PETase $\rho$ & \texttt{super} $\rho$ & Same sign \\
    \midrule
    acyl      & His catalytic-N pLDDT  & $+0.600$ & $-0.147$ & no \\
    acyl      & active-site pLDDT      & $+0.587$ & $-0.076$ & no \\
    acyl      & His pLDDT              & $+0.573$ & $-0.080$ & no \\
    acyl      & Ser pLDDT              & $+0.542$ & $-0.054$ & no \\
    acyl      & Ser-OG pLDDT           & $+0.526$ & $-0.075$ & no \\
    acyl      & ligand pLDDT (mean)    & $+0.424$ & $-0.045$ & no \\
    substrate & Ser pLDDT              & $+0.385$ & $-0.179$ & no \\
    delta     & interface PAE (min)    & $+0.371$ & $-0.177$ & no \\
    substrate & pTM                    & $+0.348$ & $-0.116$ & no \\
    delta     & His catalytic-N pLDDT  & $+0.345$ & $-0.049$ & no \\
    \midrule
    \multicolumn{2}{@{}l}{mean of these ten}       & $+0.470$ & $-0.100$ & 0/10 \\
    \multicolumn{2}{@{}l}{all 49 shared concepts}  & \multicolumn{2}{c}{$\rho_{\rho} = -0.604$} & 10/49 \\
    \bottomrule
  \end{tabular}
  \caption{\textbf{The zero-shot correlation does not survive an enzyme swap.}
  Every one of PETase's ten most strongly positive AF3 predictors inverts on a
  second serine hydrolase. Ranked by $|\rho|$ instead, one of the top ten keeps
  its sign --- the claim is about the predictors anyone would actually pick,
  which are the positive ones. $\rho_{\rho}$ is the Spearman between the two enzymes' per-metric
  Spearmans across all 49 shared concepts: $+1$ would mean a predictor can be
  picked once and reused, $0$ that one enzyme says nothing about the other, and
  a negative value that it actively misleads. Sign agreement of 10/49 is worse
  than the ${\sim}24$ expected by chance. Source:
  \texttt{figures/data/petase\_super\_af3\_transfer.csv}.}
  \label{tab:petase-super-transfer}
\end{table}

Rosetta fails differently and more mildly: $\rho_{\rho} = +0.035$ over 66 shared
concepts with 33/66 sign agreement, which is exactly chance. That is \emph{no}
transferable signal rather than an inverted one, a weaker claim than AF3's.

The nuance is worth stating, because it is the one place anything survives.
PETase's three strongest Rosetta predictors are all His--Asp distance terms, and
all three \emph{keep their sign} on \texttt{super} while collapsing in
magnitude: $-0.502 \rightarrow -0.072$, $-0.484 \rightarrow -0.023$,
$-0.410 \rightarrow -0.030$, a 7--21$\times$ shrink. Catalytic-geometry strain
points the same way on both enzymes and nearly stops mattering; it is the only
concept family here that does. Rank agreement over all 66 concepts is still nil,
so this does not rescue selection --- but it is a different failure from AF3's
outright inversion, and the direction is not random.

What \emph{does} replicate is a methodological fact about the two predictor
families. On both enzymes, Rosetta carries substantial information AF3 cannot
reconstruct while the reverse is largely false: predicting Rosetta features from
AF3 gives a median cross-validated $R^2$ of $0.145$ on PETase and $0.199$ on
\texttt{super}, against $0.556$ and $0.527$ in the other direction. The most
private features are relax-ensemble \emph{variance} terms, for which a single
structure prediction has no analogue. Same asymmetry, same magnitude, both
enzymes.

\textbf{This is the wet-lab counterpart of the transfer gap in
Section~\ref{sec:res:frontier}.} There, a policy trained inside a structure
predictor scored higher under that predictor than under a held-out one. Here, a
predictor that ranks function on the target it was calibrated against
anti-ranks it on the next target. Both say the same thing: a structure
predictor's agreement with the quantity you care about is a property of the
specific system, re-earned per target, never inherited.

That conclusion is an argument \emph{for} the library, not against it. If
zero-shot oracle scores could be trusted across targets, the right strategy
would be to rank a large virtual set and order the top of it. Because they
cannot, the ordering has to be decided by something the oracle does not control
--- which is what a combinatorially compatible library at a $67\%$ register rate
buys: enough physically assemblable, measurable diversity that the experiment,
not the predictor, does the selecting.

